\begin{corollary}[\texorpdfstring{$u=-1$}{u=-1} 处的 \texorpdfstring{$p$}{p}-进导数残差：由 Dwork 同余诱导的 \texorpdfstring{$\ZZ_p$}{Zp}-值指纹族]\label{cor:witt-minus1-padic-derivative-residues}
\leanverified{paper\_witt\_minus1\_padic\_derivative\_residues}
沿用推论 \ref{cor:witt-dwork-congruence} 的记号，取奇素数 \(p\) 并令 \(a_n(u)\in\ZZ[u]\) 为迹 ghost 序列。
记 Euler 导数算子 \(D:=u\,\frac{\mathrm{d}}{\mathrm{d}u}\)。
则对每个固定阶 \(r\ge 1\)，当 \(k>r\) 时有整除性
\[
\boxed{\ p^{\,k-r}\ \big|\ (D^r a_{p^k})(-1)\ },
\]
从而可定义整数序列
\[
R_{p,r}(k):=p^{-(k-r)}(D^r a_{p^k})(-1)\in\ZZ\qquad(k>r).
\]
并且该序列在 \(\bmod p\) 上永久稳定：
\[
\boxed{\ R_{p,r}(k)\equiv R_{p,r}(k-1)\pmod p\qquad(k>r+1)\ }.
\]
因此 \(\{R_{p,r}(k)\}_{k>r}\) 在 \(\ZZ_p\) 中为 Cauchy 列并收敛到极限
\[
R_{p,r}^{(\infty)}:=\lim_{k\to\infty}R_{p,r}(k)\in\ZZ_p,
\]
从而对每个奇素数 \(p\) 得到一族可逐级复核的 \(\ZZ_p\)-值导数残差不变量 \(\{R_{p,r}^{(\infty)}\}_{r\ge 1}\)。
\end{corollary}
\begin{proof}
由推论 \ref{cor:witt-dwork-congruence}，在 \(\ZZ[u]\) 中有
\(
a_{p^k}(u)\equiv a_{p^{k-1}}(u^p)\ (\mathrm{mod}\ p^k)
\)。
对同余施加 \(D^r\)（\(D\) 为 \(\ZZ\)-线性微分算子）仍保持模 \(p^k\)；并且对任意 \(f\) 有恒等式
\(
D(f(u^p))=p\,(Df)(u^p)
\)，
从而 \(D^r(f(u^p))=p^r(D^rf)(u^p)\)。
取 \(f=a_{p^{k-1}}\) 并在 \(u=-1\) 处特化（因 \(p\) 为奇数故 \((-1)^p=-1\)），得到
\[
(D^ra_{p^k})(-1)\equiv p^r(D^ra_{p^{k-1}})(-1)\pmod{p^k}.
\]
于是存在 \(t\in\ZZ\) 使
\(
(D^ra_{p^k})(-1)=p^r(D^ra_{p^{k-1}})(-1)+p^k t
\)。
当 \(k=r+1\) 时右端含因子 \(p\)，故 \(p\mid (D^ra_{p^{r+1}})(-1)\)；再对 \(k\) 归纳即得 \(p^{k-r}\mid (D^ra_{p^k})(-1)\)。
将上式除以 \(p^{k-r}\) 得
\[
R_{p,r}(k)=R_{p,r}(k-1)+p^r t.
\]
因 \(r\ge 1\) 有 \(p^r t\equiv 0\ (\mathrm{mod}\ p)\)，从而得到 \(\bmod p\) 的稳定性。由稳定性知其为 \(\ZZ_p\) 中的 Cauchy 列并收敛于某个 \(R_{p,r}^{(\infty)}\in\ZZ_p\)。证毕。
\end{proof}
